\documentclass[../ap_2013.tex]{subfiles}

\graphicspath{{./image/}}

\begin{document}

\setcounter{chapter}{2}
\chapter{光学:フィルム}
\section{}
スネルの法則より全反射するには
\begin{equation}\begin{aligned}[b]
    n_2\sin\theta &> n_1\\
    k_0n_2\sin\theta = \beta &> k_0n_1
\end{aligned}\end{equation}
であればよい。
また、光が\(+z\)方向に進んでいるため、
\begin{equation}\begin{aligned}[b]
    k_0n_2\sin\theta =\beta < k_0n_2
\end{aligned}\end{equation}
よって
\begin{equation}\begin{aligned}[b]
    k_0n_1< \beta < k_0n_2
\end{aligned}\end{equation}

\section{}
アンペールの式より
\begin{equation}\begin{aligned}[b]
    \begin{pmatrix}
        \partial_y H_z(x)e^{i(\omega t-\beta z)} -\partial_z H_y(x)e^{i(\omega t-\beta z)}\\
        \partial_z H_x(x)e^{i(\omega t-\beta z)} -\partial_x H_z(x)e^{i(\omega t-\beta z)}\\
        \partial_x H_y(x)e^{i(\omega t-\beta z)} -\partial_y H_x(x)e^{i(\omega t-\beta z)}
    \end{pmatrix}
    &=\epsilon_j\begin{pmatrix}
        0\\
        \partial_t E_y(x)e^{i(\omega t-\beta z)}\\
        0
    \end{pmatrix}\\
    \begin{pmatrix}
        i\beta H_y(x)\\
        -i\beta H_x(x) -\qty(\partial_x H_z(x))\\
        \qty(\partial_x H_y(x))
    \end{pmatrix}e^{i(\omega t-\beta z)}
    &=\begin{pmatrix}
        0\\
        i\omega \epsilon_j E_y(x)\\
        0
    \end{pmatrix}e^{i(\omega t-\beta z)}
\end{aligned}\end{equation}
となるので
\begin{equation}\begin{aligned}[b]
    \begin{cases}
        H_y(x) &= 0\\
    -i\beta H_x(x) -\partial_x H_z(x) &= i\omega \epsilon_j E_y(x)
    \end{cases}
\end{aligned}\end{equation}
どうようにファラデーの式でもやると
\begin{equation}\begin{aligned}[b]
    \begin{pmatrix}
        -\partial_z E_y(x)e^{i(\omega t-\beta z)}\\
        0\\
        \partial_x E_y(x)e^{i(\omega t-\beta z)}
    \end{pmatrix}
    &=-\mu_0\begin{pmatrix}
        \partial_t H_x(x)e^{i(\omega t-\beta z)}\\
        0\\
        \partial_t H_z(x)e^{i(\omega t-\beta z)}
    \end{pmatrix}\\
    \begin{pmatrix}
        i\beta E_y(x)\\
        0\\
        \partial_x E_y(x)
    \end{pmatrix}e^{i(\omega t-\beta z)}
    &=\begin{pmatrix}
        -i\omega\mu_0 H_x(x)\\
        0\\
        -i\omega\mu_0 H_z(x)
    \end{pmatrix}e^{i(\omega t-\beta z)}
\end{aligned}\end{equation}
となるので、
\begin{equation}\begin{aligned}[b]
    \begin{cases}
        \beta E_y(x) &= -\omega \mu_0H_x(x)\\
        \partial_x E_y(x) &= -i\omega \mu_0H_z(x)
    \end{cases}
\end{aligned}\end{equation}

\section{}
\(H_z(x)=0\)であるとすると、
前問で得られた方程式は
\begin{equation}\begin{aligned}[b]
    &\begin{cases}
        -i\beta H_x(x) &= i\omega\epsilon_jE_y(x)\\
        \beta E_y(x) &= -\omega\mu_0H_x(x)
    \end{cases} &\rightarrow
    &\begin{cases}
        \omega \epsilon_j E_y(x) + \beta H_x(x) &=0\\
        \beta E_y(x) + \omega \mu_0 H_x(x) &= 0
    \end{cases}
\end{aligned}\end{equation}
いま、領域内で電磁場が存在することから\(E_y\neq0\)であるため、
この\(E_y(x),H_x(x)\)の連立方程式が非自明な解を持たなければならない。
なのでその条件は
\begin{equation}\begin{aligned}[b]
    0&=\omega^2 \epsilon_2 \mu_0 -\beta^2 \\
    &= \frac{\omega^2 n_2^2}{c^2}-\beta^2\\
    &= (k_0n_2)^2-\beta^2
\end{aligned}\end{equation}
しかし、\(\beta\neq k_0n_2\)であるためこの条件を満たさず矛盾する。
よって\(H_z(x)\neq 0\)とわかる。

\section{}
ファラデーの式から変形していく。
\begin{equation}\begin{aligned}[b]
    \beta E_y(x) &= -\omega\mu_0\qty[\frac{1}{-i\beta}\qty(\partial_xH_z(x)+i\omega\epsilon_jE_y(x))]\\
    \beta^2 E_y(x) &= -i\omega\mu_0\partial_x\qty[\frac{1}{-i\omega\mu_0}\partial_xE_y(x)]
    +\omega^2\mu_0\epsilon_j E_y(x)\\
    0 &= \dv[2]{E_y(x)}{x}+\qty(k_0^2n_j^2-\beta^2)E_y(x)
\end{aligned}\end{equation}

\section{}
領域I,III では\(k_0n_1^2-\beta^2 <0\)で、
領域IIでは\(k_0n_2^2-\beta^2 >0\)であり、
\(E_y(x)\)が偶関数で無限遠で電場はなくなるので、
\begin{equation}\begin{aligned}[b]
    E_y(x) =
    \begin{cases}
        A_0\exp\qty(-\sqrt{\beta^2-k_0^2n_1^2}\abs{x})& (\abs{x}>d)\\
        A_1\cos(\sqrt{k_0^2n_2^2-\beta^2}x) & (\abs{x}<d)
    \end{cases}
\end{aligned}\end{equation}
接続条件より
\begin{equation}\begin{aligned}[b]
    \begin{cases}
        A_0\exp\qty(-\sqrt{\beta^2-k_0^2n_1^2}d) &= A_1\cos(\sqrt{k_0^2n_2^2-\beta^2}d)\\
        -\sqrt{\beta^2-k_0^2n_1^2}A_0\exp\qty(-\sqrt{\beta^2-k_0^2n_1^2}d) &= \sqrt{k_0^2n_2^2-\beta^2}A_1\sin(\sqrt{k_0^2n_2^2-\beta^2}d)
    \end{cases}
\end{aligned}\end{equation}
これの比を取って
\begin{equation}\begin{aligned}[b]
    \sqrt{k_0^2n_2^2-\beta^2}\tan(\sqrt{k_0^2n_2^2-\beta^2}d) &=\sqrt{\beta^2-k_0^2n_1^2}\\
    \tan(\sqrt{k_0^2n_2^2-\beta^2}d) &= \sqrt{\frac{\beta^2-k_0^2n_1^2}{k_0^2n_2^2-\beta^2}}
\end{aligned}\end{equation}
なので
\begin{equation}\begin{aligned}[b]
    A(\beta) &= \sqrt{k_0^2n_2^2-\beta^2}d\\
    B(\beta) &= \sqrt{\frac{\beta^2-k_0^2n_1^2}{k_0^2n_2^2-\beta^2}}
\end{aligned}\end{equation}

\section{}
\begin{equation}\begin{aligned}[b]
    F(\beta) = \sqrt{k_0^2n_2^2-\beta^2}d -\arctan(\sqrt{\frac{\beta^2-k_0^2n_1^2}{k_0^2n_2^2-\beta^2}}) = m\pi
\end{aligned}\end{equation}

\section*{感想}



\end{document}
